\documentclass[10pt,twocolumn]{article}
% NOTE: the official ICML template (icml20xx.sty) was not available offline in this
% environment, so this file mimics its layout with the article class. Swap in the
% official style file before submission.
\usepackage[margin=0.75in]{geometry}
\usepackage{booktabs,amsmath,amssymb,graphicx,microtype}
\usepackage[hidelinks]{hyperref}
\title{Predicting Random-Walk Cover Times on Small Graphs from Cheap Statistics}
\author{Autonomous research agent (unattended run)}
\date{}
\begin{document}
\maketitle

\begin{abstract}
The cover time of a simple random walk on a connected graph is the expected number of steps needed to visit every vertex. The classical bound $2m(n-1)$ is tight only up to large factors on small graphs. We ask whether a log-linear model built from four cheap statistics (vertex count $n$, edge count $m$, maximum degree, diameter) predicts the cover time of connected graphs with $n\le 30$ better than this bound, and where it fails. On 176 synthetic graphs from 11 families, measured by 2000 Monte Carlo walks per graph, the four-feature model has a median leave-one-family-out error of $\times1.26$ (95\% bootstrap CI on the median: $[0.193,0.283]$ in log units), versus $\times4.8$ for the bound. The error is concentrated in families whose cover time is set by a bottleneck (lollipops, complete graphs, binary trees), where global statistics are not informative. This is a small, single-setting smoke-test study, and the conclusions are limited to the families we generated.
\end{abstract}

\section{Introduction}
For a graph $G=(V,E)$ with $n=|V|$ and $m=|E|$, let $C(G)$ denote the expected cover time from a fixed start vertex. Aleliunas et al.\ (1979) showed $C(G)\le 2m(n-1)$, and Feige (1995) gave tight bounds in terms of $n$ alone, with worst cases attained by lollipop-like graphs. Such bounds are worst-case over all graphs, so on a specific small graph they can be off by one or more orders of magnitude. Practitioners who need a number for a given small graph (e.g.\ to size a walk-based search) would benefit from an empirical predictor that uses cheap statistics and reports its own error. We build one and measure where it fails. We do not claim a theorem.

\section{Method}
\textbf{Graphs.} We generate connected graphs with $6\le n\le 30$ from 11 families: path, cycle, star, complete, wheel, lollipop (clique plus path), grid, binary tree, random labelled tree, $G(n,p)$ conditioned on connectivity ($p\sim U[0.15,0.6]$), and random 3-regular graphs. Random families use two seeds per size. Binary trees have $2^{k}-1$ vertices, so only a subset of sizes appears. The final set has 176 graphs (see \texttt{results/cover.csv}).

\textbf{Ground truth.} For each graph we simulate $R=2000$ independent walks from vertex $0$ and record the first time all vertices are visited. The target is $y=\log \bar{T}$, where $\bar{T}$ is the sample mean. The median relative Monte Carlo standard error is $1.1\%$ (max $2.3\%$), so target noise is far below the prediction errors reported below. As a sanity check, the simulator reproduces known values: path $P_6$ from an endpoint gives $25.1$ (exact $(n-1)^2=25$), $C_{30}$ gives $438\pm 6$ (exact $n(n-1)/2=435$), and $K_{30}$ gives $114.8\pm0.8$ (exact $(n-1)H_{n-1}\approx115.9$).

\textbf{Predictors.} All predictors are ordinary least squares on $\log$ scale with an intercept:
(a) $\log n$; (b) $\log n,\log m$; (c) $\log n,\log m,\log \Delta,\log \mathrm{diam}$. We also evaluate the bound $\log(2m(n-1))$ directly, with no fitting.

\textbf{Evaluation.} The main metric is $|\log \hat{C}-\log \bar{T}|$, reported as median and 90th percentile, and as a multiplicative factor $e^{\text{median}}$. We use leave-one-family-out (LOFO) cross-validation, which tests transfer to an unseen graph family. Its 95\% CI comes from 2000 bootstrap resamples of graphs.

\section{Experiments}
Table~\ref{tab:main} summarises the predictors. The bound overestimates by a median factor of $4.8$ (range $1.7$ to $220$). Adding $\log m$ to $\log n$ lowers the LOFO error from $0.55$ to $0.41$, and adding maximum degree and diameter reduces it further, to a median factor of $1.26$. In-sample errors are only slightly smaller, so the four-feature model is not severely overfit with five parameters on 176 graphs. Its coefficients, fit on all graphs, are $\hat\beta=(-0.41,\,0.69,\,0.03,\,0.74,\,1.27)$ for (intercept, $\log n$, $\log m$, $\log\Delta$, $\log\mathrm{diam}$). The $\log m$ and $\log n$ terms partly cancel, so their individual values should not be read causally.

\begin{table}[t]
\centering\small
\caption{Prediction error in log units over 176 graphs. Median and 90th percentile of $|\log\hat C-\log\bar T|$; LOFO = leave-one-family-out.}
\label{tab:main}
\begin{tabular}{lcccc}
\toprule
Predictor & Median & p90 & Median factor & Mode \\
\midrule
$2m(n-1)$ bound & 1.57 & 3.40 & $\times4.8$ & no fit \\
$\log n$ & 0.55 & 1.04 & $\times1.7$ & LOFO \\
$\log n,\log m$ & 0.41 & 0.79 & $\times1.5$ & LOFO \\
$\log n,\log m,\log\Delta,\log\mathrm{diam}$ & 0.23 & 0.60 & $\times1.26$ & LOFO \\
\quad (same, in-sample) & 0.15 & 0.46 & $\times1.17$ & fit \\
\bottomrule
\end{tabular}
\end{table}

\begin{table}[t]
\centering\small
\caption{Per-family LOFO errors for the four-feature model. Bias is the mean signed error $\log\bar T-\log\hat C$ (positive means underprediction). Bound/true is the median ratio $2m(n-1)/\bar T$.}
\label{tab:fam}
\input{../results/table_families.tex}
\end{table}

Table~\ref{tab:fam} shows where the predictor fails. Lollipops have the largest error ($\times2.7$ median), underpredicting by a factor of about $2.5$ on average. Their cover time is governed by hitting time across the path bottleneck, which depends on the position of the clique and the length of the path rather than on $n$, $m$, $\Delta$, or diameter alone. Complete graphs and binary trees are also underpredicted (we did not investigate the binary-tree case further). Complete graphs have diameter 1, an extreme value on the log scale, so the diameter coefficient extrapolates poorly when that family is held out. Grids and $G(n,p)$ are overpredicted, and the families with the most regular structure (cycle, random 3-regular, path, random tree) are predicted to within about $\times1.1$ to $\times1.2$. Figure~\ref{fig:pred} plots residuals against the true cover time.

\begin{figure}[t]
\centering
\includegraphics[width=0.85\columnwidth]{../results/fig_pred.pdf}
\caption{LOFO residuals $\log\bar T-\log\hat C$ against $\log\bar T$, coloured by family.}
\label{fig:pred}
\end{figure}

\section{Discussion and limitations}
Cheap global statistics give a useful prediction on regular or tree-like graphs, roughly within a factor of 1.2 to 1.5, and a much better estimate than the worst-case bound. They fail on graphs with bottlenecks, where cover time depends on local structure that $n$, $m$, $\Delta$, and diameter do not capture. A practical rule is to treat the predictor as reliable only for families resembling the training set and to check for a bottleneck (e.g.\ by computing hitting times or a conductance-type quantity) before trusting it.

Limitations: (i) the graph set is small and synthetic, with 11 families and at most 30 vertices, so family-level conclusions are tentative; (ii) cover times are from a single start vertex (vertex 0), not the worst start, which for lollipops matters a lot; (iii) LOFO error is measured on families, not on new graph distributions; (iv) the regression is fit to log cover time, so predictions for a new graph with $n\le 30$ are interpolations of a small sample; (v) the model is not compared to richer predictors such as exact hitting-time computation, which is cheap at this size and would be the natural next baseline.

\section{Reproducibility}
Code is in \texttt{code/}. Run \texttt{python run\_experiments.py --out ../results/cover.csv --walks 2000 --seed 1} then \texttt{python analyze.py} (Python 3, numpy, networkx, matplotlib). The run takes about 10 seconds on one CPU core. Random seed: 1 for graph generation and walks. Outputs are in \texttt{results/}.

\section*{References}
\small
Aleliunas, R., Karp, R., Lipton, R., Lov\'asz, L., Rackoff, C. Random walks, universal traversal sequences, and the complexity of maze problems. FOCS, 1979.\\
Feige, U. A tight upper bound on the cover time for random walks on graphs. Random Structures \& Algorithms, 1995.\\
Lov\'asz, L. Random walks on graphs: a survey. Combinatorics, Paul Erd\H{o}s is Eighty, 1993.

\end{document}
